\section{Background and related work}
\label{sec:background}

\subsection{Subnational microsimulation and spatial reweighting}

Subnational microsimulation starts from a recurring problem: a national survey carries the
behavioural and demographic detail needed for policy modelling, but it does not contain enough
observations in every local area to support direct estimation. Spatial microsimulation addresses
that gap either by reweighting existing survey records or by building synthetic small-area
populations from aggregate constraints \citep{tanton2014review, odonoghue2014review}.

This distinction matters for the present work, because our pipeline is a reweighting system built on
observed CPS households. It places households across geographies to create local variants, but the
core problem remains the calibration of survey-based microdata to administrative targets, which keeps
calibration methods as the closest reference point.

The spatial microsimulation literature supplies further context. \citet{williamson1998},
\citet{huang2001}, and \citet{harland2012} describe combinatorial and synthetic-population
approaches that search over record assignments area by area. \citet{tanton2011} and
\citet{lovelace2016} discuss reweighting for small-area estimation and policy analysis. These
methods usually target one geographic layer at a time, whereas our setting asks one weighted dataset
to reproduce state and national targets together.

\subsection{Survey calibration}

Let $i = 1, \ldots, n$ index sampled units, let $d_i$ denote the initial design weight for unit $i$,
and let $\mathbf{x}_i \in \mathbb{R}^p$ denote the auxiliary variables used for calibration. Survey
calibration seeks adjusted weights $w_i$ whose weighted totals match known population totals:
\begin{equation}
    \sum_{i=1}^{n} w_i \mathbf{x}_i = \mathbf{T},
    \label{eq:calibration_constraint}
\end{equation}
where $\mathbf{T} \in \mathbb{R}^p$ is the vector of published totals \citep{deville1992,
sarndal2007}. The auxiliary variables may be counts, indicators, or continuous quantities such as
income components.

The subnational problem extends this setup in two ways. The target vector draws on several
administrative sources, so it mixes count and dollar constraints. The constraints are also nested by
geography: district totals should sum to state totals, which should sum to national totals after
uprating and reconciliation. At production scale this yields a large target system with substantial
collinearity across rows.

\subsection{Calibration estimators}

Two classical methods are the standard references for this kind of problem. Generalized regression
(GREG) minimizes a distance from the initial weights subject to the calibration equations
\citep{deville1992, sarndal2007}. It accepts quantitative targets directly and is well understood,
but it can return negative weights and becomes harder to solve as the target system grows and its
rows become collinear. Negative weights are acceptable in some estimation settings and awkward in a
microdata file meant for tax-benefit simulation, where a household is expected to represent positive
population mass.

Iterative proportional fitting \citep[IPF, or raking;][]{deming1940, ireland1968} adjusts weights
multiplicatively so that selected margins match published totals. It preserves non-negativity and
avoids matrix inversion, which makes it a natural choice for count-style margins. In spatial
microsimulation and population synthesis, raking is commonly paired with a sampling step in one of two
orders: draw a sample first and rake that fixed sample, or rake the full file first and integerise the
fractional weights into a finite population. Iterative proportional updating and related extensions
handle household and person margins simultaneously \citep{pritchard2012}. These raking-based workflows
are natural comparators for pruning methods on categorical margin surfaces. They are not direct
full-surface baselines here: classical IPF is built for non-negative categorical margins, whereas the
production target surface mixes counts, dollar totals, overlapping families, near-zero targets, and
hierarchical constraints. Generalized raking or entropy calibration can extend the idea to broader
linear constraints, but then the method becomes a general calibration optimizer rather than the
classical IPF workflow.

A third approach treats calibration as an optimization problem to be solved numerically. The two
methods above are special cases of one problem: \citet{deville1992} define the calibrated weights as
those closest to the design weights, under a chosen distance, subject to reproducing the targets,
with GREG and raking the closed-form solutions for particular distances. When the target system is
large, collinear, mixes count and dollar constraints, and restricts weights to be positive, the
classical Newton and Lagrange solvers become hard to apply or fail to converge, which motivates
minimizing the calibration objective directly. \citet{espuny2018} do so with a global optimization
that enforces the range restrictions on the weights explicitly, and minimum-distance reweighting has
long been used this way in tax microsimulation \citep{creedy2003}. At the scale of an enhanced
microdata file the practical solver is stochastic gradient descent: the weights are parameterized to
stay positive, typically in log space; a loss measures the relative discrepancy between the weighted
estimates and the targets; and an optimizer such as Adam \citep{kingma2015} minimizes it, which
scales to thousands of mixed, hierarchical targets without inverting a matrix. \citet{woodruff2024}
calibrate PolicyEngine's enhanced CPS in exactly this way.

These estimators answer a calibration question rather than a sampling one: given a fixed set of
records, they find weights that match the targets. The method studied here selects records and fits
weights at the same time, extending the gradient-based calibration just described with explicit
record selection. Classical calibration still matters for the comparison set, because it can be
combined with selection: sampling followed by raking is the classical reduce-then-calibrate analogue
to random sampling followed by gradient reweighting, while raking followed by integerisation is the
classical calibrate-then-reduce analogue to sampling from dense calibrated weights.

\subsection{Reducing microdata by selecting records}
\label{sec:reduction_methods}

Reducing a large candidate dataset to a fixed budget is a sampling problem with a long statistical
history, and several established methods reduce a microdata file by selecting actual records to match
known totals. They differ mainly in where calibration enters: before selection, after selection, or
inside the selection step itself. We review the broad option space because it defines what a reasonable
analyst might try: random sampling followed by calibration, calibration followed by integerisation,
raking-based variants of both, balanced sampling, combinatorial optimisation, and sparse optimization
relaxations. Not all of these are equally applicable to the production problem. The target surface here
is large, mixed, and hierarchical, so methods built for categorical margins or discrete area-by-area
searches are reference points rather than direct full-surface baselines. The matched-budget comparison
in Section~\ref{sec:methodology} therefore focuses on the tractable full-surface set: informed $L_0$,
informed $L_0$ followed by ordinary refitting on the selected records, random sampling with reweighting,
and survey-weight sampling. Raking-based reductions, convex sparse calibration, balanced sampling, and
combinatorial optimisation remain published comparators that motivate robustness checks where their
assumptions fit cleanly. We also note the coreset view from
machine learning, where a weighted subset is chosen so that a model fit on the subset approximates the
fit on the full data \citep{feldman2020coresets, bachem2017coresets}.

\subsubsection{Random sampling}

The simplest reduction draws records at random. Under simple random sampling without replacement each
record has inclusion probability $n/N$ and design weight $N/n$, and the Horvitz--Thompson estimator
recovers population totals without bias \citep{horvitz1952, cochran1977, sarndal1992}. A random
subsample preserves the conditional distribution of the data in expectation, but it does not
reproduce administrative totals, so it has to be calibrated after the draw. Random selection is also
the reference that the data-reduction literature treats as the genuine test: informativeness-based
schemes are measured against it \citep{ma2015leveraging}, and pruning methods in deep learning often
fail to beat it at high reduction ratios \citep{guo2022deepcore}. In our setting it is the
reduce-first, calibrate-after baseline, a random draw followed by gradient-descent reweighting to the
targets. Replacing the gradient reweighting step with raking gives the classical sampling-then-raking
workflow; it is directly comparable on a raking-compatible categorical margin subset, but not on the
full mixed count-and-dollar target surface without changing the problem into generalized calibration.

\subsubsection{Raking, survey-weight sampling, and integerisation}

A second approach calibrates first and reduces second. Given fractional weights from a calibration
step, survey-weight sampling selects records with probability proportional to those weights, which
turns a reweighted file into a finite, integer-weighted population that approximately preserves both
the calibrated totals and the distribution. This is unequal-probability,
probability-proportional-to-size sampling in the survey tradition \citep{hansen1943, horvitz1952},
applied to microsimulation as integerisation. The upstream calibration can be IPF/raking, GREG, or a
numerical optimizer; the downstream question is how to convert fractional weights into a finite file.
\citet{lovelace2013trs} catalogue the common integerisation methods and show that
truncate-replicate-sample reproduces the targets most accurately while fixing the reduced population
to an exact size; earlier top-up schemes are due to \citet{ballas2005}. Because the selection
probabilities are the calibration weights, survey-weight sampling presupposes a calibration step, so
in our comparison it is the calibrate-then-reduce baseline, with gradient descent supplying the
weights. A raking-then-TRS variant is the corresponding classical baseline on a categorical margin
surface, but it does not naturally extend to the full mixed production target surface without replacing
IPF with a broader calibration solver.

\subsubsection{Balanced sampling}

Balanced sampling selects a fixed-size sample whose Horvitz--Thompson estimates reproduce auxiliary
totals exactly or nearly exactly. The cube method of \citet{deville2004cube} is the standard
construction: it chooses inclusion indicators while preserving prescribed balancing equations, then
uses the resulting sample weights for estimation. This makes balanced sampling a target-aware
selection method rather than a post-sampling calibration method, and therefore a relevant published
neighbour of $L_0$ selection. Its fit to the present problem is not exact. Classical balanced sampling
usually treats the auxiliary totals as design constraints with predetermined inclusion probabilities,
whereas the production calibration surface here mixes categorical counts, dollar totals, and
hierarchical relationships and then fits positive weights after selection. Still, it is an important
reference point for the idea that the sample itself, not only its weights, can be chosen to preserve
known totals.

\subsubsection{Convex sparse calibration}

A final family keeps the continuous optimization setup but replaces the direct retained-record count
with a convex sparsity surrogate. An $L_1$ penalty on fitted weights, implemented with proximal
soft-thresholding, can drive weights exactly to zero while preserving the same calibration loss. This
is tractable on the full target surface and gives a useful ablation: it asks whether the value comes
from sparse joint selection in general or from the Hard Concrete $L_0$ relaxation specifically. Its
drawback is interpretive. The $L_1$ penalty shrinks weight magnitudes as well as encouraging zeros, so
its tuning parameter is not a direct retained-record-count control in the way the expected $L_0$ norm
is. We therefore treat it as a convex sparse analogue rather than as the primary sampling method.

\subsubsection{Informed selection}

These methods share the idea that a carefully chosen weighted subset can stand in for a larger
dataset, but they reach it differently: random sampling ignores the targets and repairs the gap
afterwards, survey-weight sampling inherits the targets from a prior calibration, balanced sampling
tries to preserve auxiliary totals through the sample design, and convex sparse calibration remains in
the same optimizer but changes the penalty. A fourth tradition makes selection itself the calibration:
combinatorial optimisation chooses a combination of records, with integer multiplicities, to match the
benchmark totals directly, searched with metaheuristics such as simulated annealing
\citep{kirkpatrick1983, williamson1998, voas2000}. It is the closest precedent for target-informed
selection, but it is discrete and gradient-free; we note it as related work rather than include it in
the matched-budget comparison. The method studied here also pursues the targets at selection, but with
a continuous, gradient-based relaxation that fits the selection and the weights jointly. The next
subsection sets out that relaxation.

\subsection{\texorpdfstring{$L_0$}{L0} regularization and the Hard Concrete distribution}

The $L_0$ norm counts the non-zero entries in a vector. Applied to record weights, it counts the
records that remain active, which is exactly the size of the retained sample. Selecting a subset of
a chosen size is therefore an $L_0$-penalized problem. The difficulty is that the $L_0$ norm is
non-differentiable and combinatorial, so it cannot be minimized by gradient descent in its raw form.

\citet{louizos2018} make the penalty differentiable by attaching a stochastic gate $z_i \in [0, 1]$
to each parameter and optimizing the expected $L_0$ norm. The gate is drawn from a Hard Concrete
distribution: a binary concrete (continuous-relaxation) variable stretched beyond the unit interval
and then clipped back into it. With a learned location parameter $\log \alpha_i$ and a fixed
temperature $\beta$, draw $u \sim \mathrm{Uniform}(0, 1)$ and set
\begin{align}
    s_i      &= \sigma\!\left(\frac{\log u - \log(1 - u) + \log \alpha_i}{\beta}\right), \\
    \bar{s}_i &= s_i (\zeta - \gamma) + \gamma, \\
    z_i      &= \min\!\big(1, \max(0, \bar{s}_i)\big),
    \label{eq:hardconcrete}
\end{align}
where $\sigma$ is the logistic function and $\gamma < 0 < 1 < \zeta$ are fixed stretch bounds.
Stretching $s_i$ from $(0, 1)$ to $(\gamma, \zeta)$ and clipping to $[0, 1]$ places point masses at
exactly $0$ and exactly $1$, so a gate can switch a record fully off or fully on while staying
differentiable in between.

Because the gate has a closed-form probability of being non-zero, the expected $L_0$ norm is
differentiable. The probability that gate $i$ is active is
\begin{equation}
    \Pr(z_i \neq 0) = \sigma\!\left(\log \alpha_i - \beta \log \frac{-\gamma}{\zeta}\right),
    \label{eq:l0penalty}
\end{equation}
and the $L_0$ penalty is the sum of these probabilities over all records,
$\sum_i \Pr(z_i \neq 0)$, which is the expected number of retained records; minimizing it pushes
gates toward zero. At test time the noise is removed and the gate is evaluated deterministically,
\begin{equation}
    \hat{z}_i
    =
    \min\!\left(
        1,
        \max\!\left(
            0,\,
            \sigma(\log \alpha_i)(\zeta - \gamma) + \gamma
        \right)
    \right),
    \label{eq:deterministic_gate}
\end{equation}
which returns an exact zero for any record driven far enough off and an exact one for any record driven
far enough on, with intermediate values in between. The exact zeros make the weight vector exactly
sparse; a record left with an intermediate gate is retained with its fitted weight scaled by the gate.

The two pieces depend on each other, which is the point worth stressing for readers new to this
construction. The $L_0$ penalty states the objective, that fewer active records are preferred, but
on its own it is not optimizable by gradient methods. The Hard Concrete gate makes that objective
differentiable, but on its own, without the penalty, it carries no pressure toward sparsity and
leaves every gate open. Used together they turn record selection into a term the optimizer can
minimize alongside calibration error. This is what makes the retained-record count a control
variable rather than the result of a separate thresholding step.

\subsection{\texorpdfstring{Convex sparse selection ($L_1$)}{Convex sparse selection (L1)}}
\label{sec:l1_background}

The $L_0$ count is the natural sparsity objective but a hard one. The standard convex surrogate replaces
it with the $L_1$ norm of the weights, $\lVert \mathbf{w} \rVert_1 = \sum_i |w_i|$, which is the LASSO of
\citet{tibshirani1996}. Adding the penalty to the calibration loss gives the convex program
\begin{equation}
    \min_{\mathbf{w} \ge 0}\; \mathcal{L}_{\mathrm{cal}}(\mathbf{w}) + \lambda_{L_1} \sum_{i=1}^{n} w_i,
    \label{eq:l1}
\end{equation}
where $\lambda_{L_1}$ sets the strength of the penalty and the weights are non-negative here, so
$|w_i| = w_i$. The $L_1$ norm is the standard tool for sparse regression and feature selection: unlike the
smooth $L_2$ (ridge) penalty, which shrinks coefficients but rarely sets them to zero, its corner at the
origin drives small coefficients to exact zeros. The program is convex and is minimized by proximal
gradient descent, each gradient step followed by a soft-thresholding step
$w_i \leftarrow \max(w_i - \tau,\, 0)$ that subtracts a fixed amount $\tau \propto \lambda_{L_1}$ from every
weight and clips at zero \citep{beck2009fista}, so a record whose pull on the targets does not clear the
threshold is dropped while the rest are kept and shrunk. The trade-off is structural: the single penalty
$\lambda_{L_1}$ controls both how many records survive and how much their weights shrink, because the same
norm that zeroes small weights also pulls the surviving ones toward zero, so forcing a small sample forces
heavy shrinkage on the records that remain. This convex arm is a natural next ablation for $L_0$: it
pursues the same target-informed sparse selection by a more stable route, and contrasting the two would
isolate what the $L_0$ gates buy by separating selection from weighting.

\subsection{Microsimulation dataset engines}

Beyond the calibration step itself, \citet{woodruff2024} embed that gradient-based reweighting in a
national pipeline that first imputes tax variables from the IRS Public Use File onto CPS records, and
that regularizes the weights with dropout rather than an explicit selection mechanism. The present
work keeps the gradient-based calibration but replaces dropout with the $L_0$ and Hard Concrete gates
of the previous subsection, so that record selection becomes explicit and the fitted weights are
exactly sparse. It also runs within \populace{}, a pipeline that
expresses dataset construction as a sequence of modular stages over a shared data structure. The data
sources, the imputation model, the geographic assignment, and the calibration method are
interchangeable components rather than fixed steps, which lets the same pipeline support different
countries and different methodological choices.
Section~\ref{sec:data} describes the pipeline and the configuration used here.
